This whitepaper has introduced the Large Reasoning and Action Model as a first-principles architecture for machine-assisted mathematical discovery and validated action selection. Its essence is the composition of a Prajna abstraction operator, a Grover-style amplitude amplification stage, and a three-tier validation layer, wrapped in a teacher-student distillation objective and a Reflexion outer loop. Each component has been formalised, each failure mode has been enumerated by inversion, and each claim has been tied to a concrete diagnostic metric.

The practical contribution of the whitepaper is the fully open-source proof-of-concept and proof-of-value workshop, which demonstrates that the architecture can be built from existing tools and that it behaves as predicted on toy benchmarks. The theoretical contribution is a convergence bound under a bounded-oracle assumption and a clean statement of the limits imposed by Grover's quadratic speedup and by undecidability.

The ambitious contribution is the staged roadmap toward the Clay Millennium Prize Problems. The roadmap explicitly acknowledges that the current architecture is insufficient for a frontier attack and names the prior stages that must be completed first. This stance preserves the motivating ambition of the programme without overclaiming on its present capability.

Future work will advance the roadmap stage by stage, beginning with mathlib lemmas and miniF2F-style benchmarks, and will contribute new mathematical ideas to the Prajna and higher-dimensional-lifting components that are known to be the bottleneck for deeper mathematics. The authors welcome collaboration from the formal methods, causal inference, quantum computing, and mechanistic interpretability communities.

The defensible verdict of the whitepaper is short and will be repeated here. Yes, in the weak sense: LRAM is a credible convergence accelerator for bounded, oracle-friendly mathematics. No, in the strong sense: LRAM is not a universal solver for open mathematics. Both statements are true at the same time, and together they define the intellectually honest home of the programme.
